\begin{musica}{As Carvoeiras}{Frederico Valério, Gustavo Sequeira}{Tuna Universitária do Porto}
\textbf{[}\ac{D}{Ai} ve\ac{D+5}{nham} \ac{G}{ver},
Venham ver as Carvoeiras,
Venham ver os olhos delas
Que maneiras têm de o\ac{Am}{lhar}. \ac{D7}{}__
Ai, venham \ac{Am}{ver} \ac{D7}{}__
Dois carvões numa bra\ac{Am}{sei}ra
Que puseram à ja\ac{D7}{ne}la
Para o vento os ate\ac{G}{ar}.\textbf{]} 2x \ac{C7 F}{}__ _

Sem a \ac{Fm}{cor} das bandeiras domingueiras,
Sem o \ac{F7}{Sol} que têm as \ac{Bbm}{chi}tas,
Ó Lisboa, sem fa\ac{Fm}{vor},
Como as \ac{G}{tuas} Carvo\ac{C7}{ei}ras são bo\ac{Fm}{ni}tas.

Olhai-as pas\ac{Eb}{san}do, gentis touti\ac{Ab}{ne}gras,
Por dentro tão \ac{C}{bran}cas, por fora tão \ac{Fm}{ne}gras,
As asas são \ac{Bbm}{an}cas, num ritmo \ac{Fm}{bran}do,
Porque não \ac{G}{pi}sam, deslizam vo\ac{C}{a}n\ac{A7}{do}. \ac{D}{}_

\nota{Refrão}

A correr muito cingida a ilharga,
Muito à larga o coração,
Tão bonitas são de ver
As varinas na descarga do carvão.

Gigas balouçando têm semelhanças
Com barcos vogando nas ondas das tranças.
Travessas, gaiatas, correndo, voando,
Os seios são ondas redondas arpando.

\nota{Refrão}

Ai, venham \ac{Cm}{ve}r, ai, venham \ac{G}{ver}!

\posicoes{C: x32010, C7: x32310, Cm: x35543, D: xx0232, D7: xx0212, D+5: xx0332,
  Eb: x68886, E: 022100, F: 133211, Fm: 133111, F7: 131211, G: 320033,
  Ab: 466544, Am: x02210}
\end{musica}
